Drittens scheint es wohl der Mühe wert, die Ursachen anzugeben, denen der rasche und ununterbrochene wirtschaftliche Aufschwung der Witwenkasse über einen Zeitraum von mehr als vierzig Jahren zuzuschreiben ist. Ich finde, dass dieser Aufschwung besonders durch das Zusammenwirken zweier Umstände begünstigt wurde: Mit dem einen hatte man 1794 überhaupt nicht gerechnet, mit dem anderen hätte man zumindest nicht rechnen dürfen.

Die erste Ursache ist der Anstieg des Zinsfußes, der schon wenige Jahre nach jenem Zeitpunkt begann. Wie ich nachgewiesen habe, glaubte man damals, mit Sicherheit nur auf 3 Prozent rechnen zu dürfen. Schon 1799 bemerkte jedoch das Kuratorium im oben bereits angeführten Reskript vom 23. Mai, dass Gelegenheiten zu einer sicheren Anlage zu 4 Prozent keineswegs selten seien, und bot sogar seine eigene Mitwirkung dabei an. Etwas später stieg der Zinsfuß allgemein auf 5 Prozent und blieb für den größten Teil der Kapitalien der Witwenkasse über eine lange Reihe von Jahren auf dieser Höhe.

Die zweite Ursache ist, dass die Zahl der Pensionen während jenes langen Zeitraums \textit{unter} dem zu erwartenden Mittelwert (15) geblieben ist. Ja, man muss sagen: \textit{beträchtlich} unter dem zu erwartenden Mittelwert, wenn man dafür nach dem plausibleren Überschlag S. [145] die Zahl 17 annimmt (nämlich \(\tfrac{4}{7}\times26\) Witwen mit einem Zuschlag von \(\tfrac{1}{6}\) wegen der Waisen). Dass man diese Erscheinung keineswegs als etwas sehr Außerordentliches betrachten soll, wurde schon oben S. [144, Z. 7] erwähnt. Ebenso wenig wäre es aber etwas Außerordentliches gewesen, wenn genau das Gegenteil eingetreten und etwa schon 20—25 Jahre nach jenem Zeitpunkt von 1794 der Mittelwert auf der anderen Seite beträchtlich und andauernd überschritten worden wäre. In einem solchen Fall wäre an die Stelle des raschen Anstiegs des Kassenvermögens ein sehr langsamer getreten, vielleicht sogar ein Stillstand oder die Notwendigkeit, auf die Quittungsklausel zurückzugreifen, falls zugleich ein starker Rückgang des Zinsfußes hinzugekommen wäre.

Eine höhere Witwenzahl ist nun schon seit einer Reihe von Jahren eingetreten, S. [126, Z. 23] und
\SourceLeaf{154}
S. [128, Z. 29]. Der Zinsfuß ist jedoch, obwohl er von seiner früheren Höhe stark zurückgegangen ist, noch immer hoch genug, dass die Kasse jener höheren Pensionenzahl gewachsen bleibt. Überhaupt hat diese hohe Zahl, selbst wenn sie noch um eine oder ein paar Pensionen mehr gestiegen wäre, an sich und insofern man darin nur ein ungewöhnlich hohes Schwanken über den Mittelwert 15 oder 17 erkennt, bei Weitem nicht eine so schwerwiegende Bedeutung wie \textit{der} Umstand, zu dem ich jetzt übergehe
\begin{quote}
dass der Mittelwert der Pensionenzahl, ob man nun 15 oder 17 für den plausibelsten hält, nur so lange gilt, wie die Genossenschaft keinen größeren Umfang erreicht als 1794, und dass diese Gültigkeit jetzt ganz aufgehört hat, da ihr Umfang, gemessen an der Zahl der verheirateten Mitglieder, um mehr als 60 Prozent größer ist als damals.
\end{quote}
\noindent\hspace*{\parindent}Für die Zeit zwischen 1794 und 1845 lässt sich dieser Maßstab nicht anwenden, weil die Zahl der verheirateten Mitglieder nicht bekannt ist. Soweit sich aber aus der Entwicklung der Gesamtzahl aller Mitglieder schließen lässt, war der Umfang der Genossenschaft im Ganzen und abgesehen von einigen Schwankungen bis etwa 1831 nicht größer, sondern eher etwas geringer als 1794. Erst seit wenigen Jahren ist die Vergrößerung \textit{sehr bedeutend} geworden. Der erweiterte Umfang der Genossenschaft steht daher mit der hohen Witwenzahl in letzter Zeit überhaupt nicht in einem ursächlichen Zusammenhang, wie bereits oben [S. 128] bemerkt wurde.

Der vierte hier noch zu betrachtende Punkt und gleichsam der Schlussstein der ersten Abteilung ist, einen Überschlag über die künftige Entwicklung der Witwenzahl im Allgemeinen zu machen, soweit sich ein solcher auf dem bisher eingeschlagenen Weg gewinnen lässt. Die Gesamtzahl der Professoren beträgt jetzt 57; davon nehmen 51 an der Witwenkasse teil, und unter diesen sind 42 verheiratet. Macht man zunächst den Überschlag nach der Hypothese, die den Verhältnissen der hiesigen Professorenwitwenkasse am meisten angemessen scheint, dass auf 7 bestehende Ehen 4 Witwen zu rechnen sind, so ergibt die Rechnung 24 Witwen. Das bedeutet, dass nach Erreichen des Beharrungszustands durch die Gesellschaft von jenem Umfang die durchschnittliche Witwenzahl 24 sein wird. Wegen der Waisenpensionen müsste nach \textsc{Kritter}s Gutachten noch ein Sechstel hinzugerechnet werden. Ich will jedoch nur ein Achtel ansetzen, die durchschnittliche Pensionenzahl im Beharrungszustand also \(=27\) setzen. Allerdings ist der Beharrungszustand erst nach 45—50 Jahren zu erwarten. Man würde sich aber sehr täuschen, wenn man diesen großen zeitlichen Abstand für einen starken Grund zur Beruhigung hielte. Denn schon lange vorher ist man dem Grenzzustand so nahe gekommen, dass der Unterschied nicht mehr viel bedeutet. Man vergleiche die oben S. [141] für einen allmählichen Anstieg der Witwenzahl angegebenen Ansätze: Ohne Anspruch auf strenge Genauigkeit können sie doch einigermaßen eine Vorstellung vom Verlauf geben. Bedenkt man dabei, dass der frühere Umfang der Gesellschaft bereits 17 Pensionen als Normalzahl (beim Verhältnis \(7:4\)) ergeben hatte, und dass es hier also nur um die allmähliche Entwicklung der auf 10 geschätzten Vergrößerung geht, so wird man mich leicht verstehen, wenn ich behaupten muss, dass man schon nach etwa 25 Jahren auf 25 Pensionen als \textit{Mittelzahl} gefasst sein müsse. Hinzu kommt nun noch die Erweiterung des Spielraums wegen der Schwankungen, denen sich unmöglich scharfe Grenzen setzen lassen. So viel ist aber gewiss: Ein Schwanken von 7 Pensionen über und unter dem Mittelwert muss als etwas keineswegs Außerordentliches in den Überschlag aufgenommen werden, da frühere Beispiele ein verhältnismäßig mindestens ebenso großes Schwanken tatsächlich belegen. Das Ergebnis dieser Überlegungen ist also, dass man um das Jahr 1870 die Pensionenzahl zwischen 18 und 32 erwarten muss, ohne im Voraus bestimmen zu können, \textit{wo} innerhalb dieses weiten Spielraums; dass das Erreichen des einen oder anderen Extrems nicht als etwas sehr Außerordentliches gelten darf; schließlich, dass diese Zahlen später noch ein
\SourceLeaf{155}
wenig vergrößert werden müssen, sodass nach längerer Zeit der Spielraum mit 20—34 zu bezeichnen ist. Das ist alles, was die Wahrscheinlichkeitsrechnung lehren kann, \textit{solange man die Gesellschaft gleichsam nur als Gesamtheit betrachtet.} Vielleicht meint mancher, das sei wenig! Ich sähe das doch anders. Es ist sehr wichtig, das Maß der Erwartungen zu kennen, zu denen man berechtigt ist, und sich nicht einer unbegründeten Sicherheit zu überlassen. Betrachten wir die beiden Extreme. Es ist \textit{möglich}, dass die Pensionenzahl nach 25 Jahren noch nicht über 18 gestiegen ist und sich auch noch geraume Zeit länger auf dieser oder einer nur wenig größeren Höhe hält. Sie kann auch vorher von ihrer jetzigen Höhe (19) zunächst noch zurückgehen und so von jetzt bis 1870 im Durchschnitt nicht unbeträchtlich unter 18 liegen. Auf diese Weise kann die Kasse Kräfte sammeln, mit denen sie, wenn sich erst nach langer Zeit das Blatt wendet, auch den extremen Fällen auf der anderen Seite die Spitze bieten kann. Das Schlimmste wäre eigentlich, wenn diese günstigen Voraussetzungen sich noch weiter verwirklichten. Ich meine: wenn in dieser Zwischenzeit die Pensionenzahl zunächst noch einmal unter 16 sänke und die Kasse dadurch nach der bestehenden Satzung noch mit einer neuen Pensionserhöhung von 10, 20, vielleicht gar 30 Rthl. belastet würde. — Wäre umgekehrt, was ebenso möglich ist, die Pensionenzahl schon nach 25 Jahren auf oder nahe an 32 gestiegen, so genügt ein grober Überschlag, um sich zu überzeugen, dass die Kasse, abgesehen von ganz außerordentlichen Zuflüssen, solchen Anforderungen nicht gewachsen sein wird, sondern, vermutlich schon früher, ihr Unvermögen wird erklären müssen.

Dass die Überschläge etwas günstiger ausfallen, wenn man statt des Verhältnisses \(7:4\) das von \(2:1\) annimmt, versteht sich von selbst. Man hätte dann nach 25 Jahren einen Spielraum von 16—28 Pensionen und nach noch längerer Zeit einen solchen von 18—30 zu erwarten.

Schließlich muss ich noch bemerken, dass dabei stillschweigend vorausgesetzt ist, der Umfang der Genossenschaft bleibe künftig auf seiner gegenwärtigen Höhe. Nimmt er weiter zu, muss man sich auf verhältnismäßig noch größere Zahlen gefasst machen. Verringert er sich dagegen wieder, darf man auch von den obigen Ergebnissen einen gewissen Abzug machen. Nach der Natur der Sache ist es aber wohl wenig wahrscheinlich, dass eine \textit{bedeutende} Verringerung anders als nach längerer Zeit eintreten könnte.

